%% kapa: exam=1 round=36 year=2025 period=1 form=A subject=civil_law subject_ko=민법 seq=2 number=2
%% source: Q-Net 감정평가사 기출문제 (cst00309) artlSeq=5231525 file=raw/1008/5231525/2025년 제36회 감정평가사 1차시험 1교시 문제지 원본.pdf
%% conversion: kordoc 4.24.51 -> tools/gen1.py (auto); review=proofread-sonnet; flags=none
\begin{question}{2}
제한능력자에 관한 설명으로 옳지 않은 것은?
\begin{choices}
  \item 미성년자가 법정대리인으로부터 허락을 얻은 특정한 영업에 관하여는 성년자와 동일한 행위능력이 있다.
  \item 가정법원은 취소할 수 없는 피성년후견인의 법률행위의 범위를 정할 수 있다.
  \item 한정후견인은 일상생활에 필요하고 그 대가가 과도하지 아니한 피한정후견인의 법률행위를 취소할 수 없다.
  \item 범위를 정하여 미성년자의 재산처분을 허락한 법정대리인은 미성년자가 그 재산을 처분하기 전에 그 허락을 취소할 수 있다.
  \item 성년후견개시의 심판은 본인의 의사에 반하여 할 수 없다.
\end{choices}
\end{question}
\answer{5}
